\section{2D Human Pose Estimation}
\label{sec:related_work_2d_hpe}
\subsection{Introduction}
2D Human Pose Estimation (HPE) has been a significant area of research in computer vision, with numerous approaches developed over the years. The 
task involves detecting and localizing human joints in images or videos, which is crucial for various applications such as activity recognition, 
human-computer interaction, and augmented reality. 2D keypoints extrapolation is often a preliminary step for 3D pose estimation, making it a 
foundational component in the broader field of human pose analysis.

The 2D pose estimation goals is to predict the coordinates of keypoints corresponding to human joints in a 2D image. With predict coordinates, we
mean the \((x, y)\) pixel locations of each joint in the image frame. The keypoints typically include joints such as the head, shoulders, elbows,
 wrists, hips, knees, and ankles.

When you are working with a 2D pose estimation model, the input is usually an image or a sequence of images (in the case of video). The output 
is a set of keypoints, containing the \((x, y)\) coordinates for each joint detected in the image, that represent the estimated pose of the person 
or people in the image.

\subsection{Evolution}
The evolution of 2D human pose estimation has seen significant advancements, particularly with the advent of deep learning techniques, I will
discuss only this type of approach in this section, that have revolutionized the field, with the advent of DeepPose.

The first tendency was to use Convolutional Neural Networks (CNNs) to directly regress the coordinates of keypoints from images. But this approach
faced challenges due to the complexity of human poses and variations in appearance, lighting, and occlusions. So researchers shifted towards 
heatmap-based methods, where the model predicts a probability distribution over possible locations for each keypoint, allowing for more robust 
localization that takes into account uncertainties in the predictions of joint positions.

TODO: DA FINIRE

\subsection{Principal approaches}
When we talk about 2D human pose estimation, we can divide the approaches into two main categories: single pose estimation and multi-pose
estimation.
In single pose estimation, the two main approaches are regression-based methods and heatmap-based methods. Regression-based methods directly 
predict the coordinates of keypoints from the input image, while heatmap-based methods generate heatmaps for each keypoint, indicating the 
likelihood of the keypoint's presence at each pixel location.
In multi-pose estimation, the two main approaches are top-down methods and bottom-up methods. Top-down methods first detect individual persons
in the image and then estimate the pose for each detected person separately. Bottom-up methods, on the other hand, first detect all keypoints 
in the image and then group them into individual poses. The choice between these approaches often depends on the specific application and the 
trade-offs between accuracy and computational efficiency, a top-down approach is generally more accurate but computationally expensive, while a
bottom-up approach is more efficient but may struggle with occlusions and overlapping poses.

In single pose video based estimation, is now common to use a single pose estimation model to extract 2D keypoints from each frame of the video,
and then thought about the temporal consistency between frames to improve the accuracy and stability of the pose estimation over time.
In multi pose video based estimation, there are two main strategies: online approaches and offline approaches. Online approaches process each 
frame sequentially, estimating poses in real-time as the video is being captured. Offline approaches, on the other hand, analyze the entire video 
after it has been recorded, allowing for more complex temporal modeling and potentially higher accuracy, as they can leverage information from 
future frames.

\subsection{Datasets, losses and metrics}




\subsection{Single Pose Estimation}


\subsection{Multi Pose Estimation}